A linguagem C foi criada em 1970 por Dennis Ritchie com a finalidade de 
desenvolver sistemas operacionais, drivers e outros softwares de baixo 
nível. Apesar de ter mais de 60 anos, continua sendo amplamente 
utilizada em diversos tipos de hardware, desde supercomputadores até 
microcontroladores e sistemas embarcados.

Com compiladores disponíveis em praticamente todos os sistemas modernos, 
destacou-se pela portabilidade e foi amplamente empregada na construção 
de utilitários e aplicações complexas que exigiam maior velocidade e 
eficiência. Em 1978, foi publicado o livro "The C Programming Language", 
escrito pelo coautor da linguagem C, que por muito tempo serviu como 
referência padrão para a linguagem.

Por ser uma linguagem imperativa, com suporte à programação estruturada, 
amplo escopo lexical e acesso eficiente à memória do computador em 
baixo nível, sendo frequentemente a escolha ideal para projetos em 
sistemas com grandes restrições de hardware.
